% GENERATED FILE. Edit canonical lessons, then run npm run generate:books.
% canonical-source-hash: fnv1a64:ab54cf505370b429
% canonical-lessons: GU-C222-prthvi, GU-C222-vatavaran, GU-C222-kudrat, GU-C222-bharti, GU-C222-ot

\chapter{Earth air and sea}
\label{ch:gu-earth-air-and-sea}
\hlchaptermodality{222}

\begin{chapteropening}
\textbf{By the end of this chapter:} \emph{I can say what the earth and the sea will do.}

Complete the last lesson of chapter 222: I can say what the earth and the sea will do.
\end{chapteropening}

\section[\texorpdfstring{\gu{પૃથ્વી} (pṛthvī)}{પૃથ્વી (pṛthvī)}]{\gu{પૃથ્વી} (pṛthvī) — planet earth, the world}

\label{lesson:GU-C222-prthvi}



\begin{quote}
Before the new one: say the Gujarati for to catch fire, then the Gujarati for earthquake.
\end{quote}

\subsection*{You'll want to know: \gu{પૃથ્વી}}
\textbf{\gu{પૃથ્વી}} — \emph{pṛthvī} — \textquotedblleft{}planet earth, the world\textquotedblright{}.

\begin{quote}
\textbf{\gu{પૃથ્વી} \gu{હંમેશા} \gu{સુંદર} \gu{રહેશે}.} — \emph{pṛthvī hammeśā sundar raheśe} — \textquotedblleft{}The earth will always stay beautiful.\textquotedblright{}
\end{quote}

A feminine noun.

\begin{cousinweb}
\textbf{\gu{પૃથ્વી}} (\emph{pṛthvī}) is Sanskrit \emph{pṛthvī}, \textquotedblleft{}the earth\textquotedblright{}, taken over whole as a learned loan, a tatsama. CDIAL lists the inherited Old Gujarati \emph{puhavī}, \textquotedblleft{}the earth\textquotedblright{}, under the same entry (CDIAL 8357).
\end{cousinweb}

\begin{grammarlens}[title={What you've built}]
One more word for everyday life at A2.
\end{grammarlens}

\subsection*{Your turn}
\begin{itemize}
\raggedright
  \item \emph{Say it:} \emph{pṛthvī}
  \item \emph{Say it:} \emph{pṛthvī}, once more
  \item \emph{Say it:} \emph{pṛthvī hammeśā sundar raheśe} — The earth will always stay beautiful.
  \item \emph{Recall:} say the Gujarati for to catch fire, then the Gujarati for earthquake, then say \emph{pṛthvī} again
\end{itemize}

\begin{tcolorbox}[breakable,colback=teal!4,colframe=teal!35!black,title={Before you move on}]
What does \gu{પૃથ્વી} mean? (\textquotedblleft{}Planet earth, the world\textquotedblright{}.) Say it once more.
\end{tcolorbox}
\section[\texorpdfstring{\gu{વાતાવરણ} (vātāvaraṇ)}{વાતાવરણ (vātāvaraṇ)}]{\gu{વાતાવરણ} (vātāvaraṇ) — atmosphere, surroundings}

\label{lesson:GU-C222-vatavaran}



\begin{quote}
Before the new one: say the Gujarati for earthquake, then the Gujarati for planet earth.
\end{quote}

\subsection*{You'll want to know: \gu{વાતાવરણ}}
\textbf{\gu{વાતાવરણ}} — \emph{vātāvaraṇ} — \textquotedblleft{}atmosphere, surroundings\textquotedblright{}.

\begin{quote}
\textbf{\gu{ગામનું} \gu{વાતાવરણ} \gu{શાંત} \gu{હશે}.} — \emph{gāmnũ vātāvaraṇ śānt haśe} — \textquotedblleft{}The atmosphere in the village will be calm.\textquotedblright{}
\end{quote}

\begin{cousinweb}
\textbf{\gu{વાતાવરણ}} (\emph{vātāvaraṇ}) is a learned word built from Sanskrit \emph{vāta}, \textquotedblleft{}wind, air\textquotedblright{}, and \emph{āvaraṇa}, \textquotedblleft{}covering\textquotedblright{} (CDIAL 1413), so the atmosphere is the air that wraps the earth; EtymDB lists it as borrowed from Sanskrit.
\end{cousinweb}

\begin{grammarlens}[title={What you've built}]
One more word for everyday life at A2.
\end{grammarlens}

\subsection*{Your turn}
\begin{itemize}
\raggedright
  \item \emph{Say it:} \emph{vātāvaraṇ}
  \item \emph{Say it:} \emph{vātāvaraṇ}, once more
  \item \emph{Say it:} \emph{gāmnũ vātāvaraṇ śānt haśe} — The atmosphere in the village will be calm.
  \item \emph{Recall:} say the Gujarati for earthquake, then the Gujarati for planet earth, then say \emph{vātāvaraṇ} again
\end{itemize}

\begin{tcolorbox}[breakable,colback=teal!4,colframe=teal!35!black,title={Before you move on}]
What does \gu{વાતાવરણ} mean? (\textquotedblleft{}Atmosphere, surroundings\textquotedblright{}.) Say it once more.
\end{tcolorbox}
\section[\texorpdfstring{\gu{કુદરત} (kudrat)}{કુદરત (kudrat)}]{\gu{કુદરત} (kudrat) — nature}

\label{lesson:GU-C222-kudrat}



\begin{quote}
Before the new one: say the Gujarati for planet earth, then the Gujarati for atmosphere.
\end{quote}

\subsection*{You'll want to know: \gu{કુદરત}}
\textbf{\gu{કુદરત}} — \emph{kudrat} — \textquotedblleft{}nature\textquotedblright{}.

\begin{quote}
\textbf{\gu{કુદરત} \gu{અમને} \gu{પાણી} \gu{આપશે}.} — \emph{kudrat amne pāṇī āpśe} — \textquotedblleft{}Nature will give us water.\textquotedblright{}
\end{quote}

A feminine noun.

\begin{cousinweb}
\textbf{\gu{કુદરત}} (\emph{kudrat}) is a loan from Arabic \emph{qudra}, \textquotedblleft{}power\textquotedblright{}, which reached Gujarati through Persian, so nature is the power that makes things. Its Arabic root is \emph{q-d-r} (Lane).
\end{cousinweb}

\begin{grammarlens}[title={What you've built}]
One more word for everyday life at A2.
\end{grammarlens}

\subsection*{Your turn}
\begin{itemize}
\raggedright
  \item \emph{Say it:} \emph{kudrat}
  \item \emph{Say it:} \emph{kudrat}, once more
  \item \emph{Say it:} \emph{kudrat amne pāṇī āpśe} — Nature will give us water.
  \item \emph{Recall:} say the Gujarati for planet earth, then the Gujarati for atmosphere, then say \emph{kudrat} again
\end{itemize}

\begin{tcolorbox}[breakable,colback=teal!4,colframe=teal!35!black,title={Before you move on}]
What does \gu{કુદરત} mean? (\textquotedblleft{}Nature\textquotedblright{}.) Say it once more.
\end{tcolorbox}
\section[\texorpdfstring{\gu{ભરતી} (bhartī)}{ભરતી (bhartī)}]{\gu{ભરતી} (bhartī) — high tide}

\label{lesson:GU-C222-bharti}



\begin{quote}
Before the new one: say the Gujarati for atmosphere, then the Gujarati for nature.
\end{quote}

\subsection*{You'll want to know: \gu{ભરતી}}
\textbf{\gu{ભરતી}} — \emph{bhartī} — \textquotedblleft{}high tide\textquotedblright{}.

\begin{quote}
\textbf{\gu{સાંજે} \gu{દરિયામાં} \gu{ભરતી} \gu{આવશે}.} — \emph{sā̃je dariyāmā̃ bhartī āvśe} — \textquotedblleft{}High tide will come in the evening.\textquotedblright{}
\end{quote}

A feminine noun.

\begin{cousinweb}
\textbf{\gu{ભરતી}} (\emph{bhartī}) is built on the inherited verb \emph{bharvũ}, \textquotedblleft{}to fill\textquotedblright{}, so high tide is the sea filling up. CDIAL derives \emph{bharvũ} from Sanskrit \emph{bharati}, \textquotedblleft{}bears, brings\textquotedblright{} (CDIAL 9397).
\end{cousinweb}

\begin{grammarlens}[title={What you've built}]
One more word for everyday life at A2.
\end{grammarlens}

\subsection*{Your turn}
\begin{itemize}
\raggedright
  \item \emph{Say it:} \emph{bhartī}
  \item \emph{Say it:} \emph{bhartī}, once more
  \item \emph{Say it:} \emph{sā̃je dariyāmā̃ bhartī āvśe} — High tide will come in the evening.
  \item \emph{Recall:} say the Gujarati for atmosphere, then the Gujarati for nature, then say \emph{bhartī} again
\end{itemize}

\begin{tcolorbox}[breakable,colback=teal!4,colframe=teal!35!black,title={Before you move on}]
What does \gu{ભરતી} mean? (\textquotedblleft{}High tide\textquotedblright{}.) Say it once more.
\end{tcolorbox}
\section[\texorpdfstring{\gu{ઓટ} (oṭ)}{ઓટ (oṭ)}]{\gu{ઓટ} (oṭ) — low tide, ebb}

\label{lesson:GU-C222-ot}



\begin{quote}
Before the new one: say the Gujarati for nature, then the Gujarati for high tide.
\end{quote}

\subsection*{You'll want to know: \gu{ઓટ}}
\textbf{\gu{ઓટ}} — \emph{oṭ} — \textquotedblleft{}low tide, ebb\textquotedblright{}.

\begin{quote}
\textbf{\gu{સવારે} \gu{ઓટ} \gu{હશે}.} — \emph{savāre oṭ haśe} — \textquotedblleft{}It will be low tide in the morning.\textquotedblright{}
\end{quote}

A feminine noun.

\begin{cousinweb}
\textbf{\gu{ઓટ}} (\emph{oṭ}) is a tadbhava. CDIAL lists it, glossed \textquotedblleft{}ebb tide\textquotedblright{}, under \emph{apavṛtti}, the noun beside Sanskrit \emph{apavartate}, \textquotedblleft{}turns back\textquotedblright{}, so the ebb is the sea turning back (CDIAL 452).
\end{cousinweb}

\begin{grammarlens}[title={What you've built}]
One more word for everyday life at A2.
\end{grammarlens}

\subsection*{Your turn}
\begin{itemize}
\raggedright
  \item \emph{Say it:} \emph{oṭ}
  \item \emph{Say it:} \emph{oṭ}, once more
  \item \emph{Say it:} \emph{savāre oṭ haśe} — It will be low tide in the morning.
  \item \emph{Recall:} say the Gujarati for nature, then the Gujarati for high tide, then say \emph{oṭ} again
\end{itemize}

\begin{tcolorbox}[breakable,colback=teal!4,colframe=teal!35!black,title={Before you move on}]
What does \gu{ઓટ} mean? (\textquotedblleft{}Low tide, ebb\textquotedblright{}.) Say it once more.
\end{tcolorbox}
